% =====================================================================
%  minggu-05.tex  --  Slide Pertemuan Minggu ke-5
%  Pengantar Kriptografi (MPFA014) -- Prodi Matematika FMIPA UBB
%
%  Topik silabus : Kriptografi klasik II
%  Sumber materi : buku_ajar/chapters/04/section-02.tex, section-03.tex,
%                  dan section-04.tex; bagian analisis frekuensi pada
%                  section-01.tex
%  Sub-CPMK      : 1.1
%
%  Cakupan: analisis frekuensi, Vigen\`ere, metode Kasiski, Playfair,
%           Hill, dan Enigma. One-time pad (minggu 6), Caesar dan cipher
%           monoalfabetik (minggu 4), serta cipher modern tidak dibahas
%           pada pertemuan ini.
%
%  Kompilasi dari dalam folder ppt/:
%      compile-ppt.bat 05
% =====================================================================

\documentclass[aspectratio=169,11pt]{beamer}
% =====================================================================
%  beamer-preamble.tex  --  prasarana dan gaya slide kuliah
%
%  Mata kuliah : Pengantar Kriptografi (MPFA014)
%  Program     : Program Studi Matematika, FMIPA, Universitas Bale Bandung
%
%  Berkas ini TIDAK memuat \documentclass. Setiap berkas minggu memuatnya
%  sendiri lalu menarik berkas ini:
%
%      \documentclass[aspectratio=169,11pt]{beamer}
%      % =====================================================================
%  beamer-preamble.tex  --  prasarana dan gaya slide kuliah
%
%  Mata kuliah : Pengantar Kriptografi (MPFA014)
%  Program     : Program Studi Matematika, FMIPA, Universitas Bale Bandung
%
%  Berkas ini TIDAK memuat \documentclass. Setiap berkas minggu memuatnya
%  sendiri lalu menarik berkas ini:
%
%      \documentclass[aspectratio=169,11pt]{beamer}
%      % =====================================================================
%  beamer-preamble.tex  --  prasarana dan gaya slide kuliah
%
%  Mata kuliah : Pengantar Kriptografi (MPFA014)
%  Program     : Program Studi Matematika, FMIPA, Universitas Bale Bandung
%
%  Berkas ini TIDAK memuat \documentclass. Setiap berkas minggu memuatnya
%  sendiri lalu menarik berkas ini:
%
%      \documentclass[aspectratio=169,11pt]{beamer}
%      \input{beamer-preamble}
%      \title{...} \subtitle{...} \setWeekInfo{Minggu n}{Topik}
%
%  Berkas minggu hanya berisi \setWeekInfo dan lingkungan frame.
%  Kompilasi dijalankan DARI DALAM folder ppt/ (lihat compile-ppt.bat),
%  karena gambar buku dirujuk dengan jalur relatif ../buku_ajar/.
% =====================================================================

% ---------------------------------------------------------------- bahasa
\usepackage[utf8]{inputenc}
\usepackage[T1]{fontenc}
\usepackage[indonesian]{babel}
\usepackage{lmodern}
\usepackage{microtype}

% ------------------------------------------------------------- matematika
\usepackage{amsmath,amssymb,mathtools}

% ----------------------------------------------------------------- tabel
\usepackage{booktabs}
\usepackage{array}
\usepackage{tabularx}

% --------------------------------------------------------------- grafik
\usepackage{tikz}
\usetikzlibrary{arrows.meta,positioning,shapes.geometric,calc,fit,backgrounds}

% ---------------------------------------------------------- kode program
\usepackage{listings}
\lstdefinestyle{slidecode}{
    basicstyle=\ttfamily\scriptsize,
    keywordstyle=\color{ubbblue}\bfseries,
    commentstyle=\color{ubbgreen}\itshape,
    stringstyle=\color{ubbred},
    numbers=left,
    numberstyle=\tiny\color{gray},
    numbersep=6pt,
    frame=single,
    rulecolor=\color{ubbblue!40},
    backgroundcolor=\color{ubbblue!3},
    breaklines=true,
    showstringspaces=false,
    tabsize=4,
    columns=fullflexible,
    keepspaces=true,
}
\lstset{style=slidecode}

% =====================================================================
%  WARNA INSTITUSI
% =====================================================================
\definecolor{ubbblue}{RGB}{12, 62, 120}      % warna utama
\definecolor{ubbteal}{RGB}{0, 118, 110}      % aksen kedua
\definecolor{ubbgreen}{RGB}{22, 122, 74}     % untuk blok contoh
\definecolor{ubbred}{RGB}{168, 42, 42}       % untuk blok peringatan
\definecolor{ubbgray}{RGB}{70, 76, 84}       % teks sekunder

% =====================================================================
%  IDENTITAS MATA KULIAH
% =====================================================================
\newcommand{\CourseName}{Pengantar Kriptografi}
\newcommand{\CourseCode}{MPFA014}
\newcommand{\CourseSemester}{Semester Ganjil}
\newcommand{\CourseProgram}{Program Studi Matematika}
\newcommand{\CourseFaculty}{Fakultas Matematika dan Ilmu Pengetahuan Alam}
\newcommand{\CourseInstitution}{Universitas Bale Bandung}
\newcommand{\CourseLecturer}{Cecep Suwanda, S.Si., M.Kom.}

% Label minggu, disetel tiap berkas lewat \setWeekInfo
\newcommand{\WeekLabel}{}
\newcommand{\WeekTopic}{}
\newcommand{\setWeekInfo}[2]{%
    \renewcommand{\WeekLabel}{#1}%
    \renewcommand{\WeekTopic}{#2}%
}

% =====================================================================
%  TEMA: header, footer, warna blok
% =====================================================================
\usetheme{default}
\setbeamertemplate{navigation symbols}{}
\setbeamertemplate{blocks}[rounded][shadow=false]

\setbeamercolor{palette primary}{bg=ubbblue, fg=white}
\setbeamercolor{palette secondary}{bg=ubbgray, fg=white}
\setbeamercolor{palette tertiary}{bg=ubbteal, fg=white}

\setbeamercolor{structure}{fg=ubbblue}
\setbeamercolor{normal text}{fg=black!92}
\setbeamercolor{alerted text}{fg=ubbred}

\setbeamercolor{block title}{bg=ubbblue, fg=white}
\setbeamercolor{block body}{bg=ubbblue!6, fg=black!92}
\setbeamercolor{block title alerted}{bg=ubbred, fg=white}
\setbeamercolor{block body alerted}{bg=ubbred!6, fg=black!92}
\setbeamercolor{block title example}{bg=ubbgreen, fg=white}
\setbeamercolor{block body example}{bg=ubbgreen!6, fg=black!92}

\setbeamercolor{itemize item}{fg=ubbblue}
\setbeamercolor{itemize subitem}{fg=ubbteal}
\setbeamercolor{enumerate item}{fg=ubbblue}

\setbeamerfont{headline}{size=\scriptsize}
\setbeamerfont{footline}{size=\scriptsize}
\setbeamerfont{block title}{size=\normalsize}

% Badan blok rata kiri (bukan rata kanan-kiri): teks slide lebih mudah dibaca
% dan baris terakhir yang pendek tidak lagi memicu "Underfull \hbox".
\addtobeamertemplate{block body begin}{\raggedright}{}
\addtobeamertemplate{block alerted body begin}{\raggedright}{}
\addtobeamertemplate{block example body begin}{\raggedright}{}

% Header: nama mata kuliah dan kode MK
\setbeamertemplate{headline}{%
    \begin{beamercolorbox}[wd=\paperwidth,ht=2.6ex,dp=1.2ex,leftskip=1em,rightskip=1em]{palette primary}%
        \usebeamerfont{headline}\CourseName\ \textperiodcentered\ \CourseCode%
        \hfill \CourseProgram%
    \end{beamercolorbox}%
}

% Footer: minggu dan nomor halaman
\setbeamertemplate{footline}{%
    \begin{beamercolorbox}[wd=\paperwidth,ht=2.4ex,dp=1.2ex,leftskip=1em,rightskip=1em]{palette secondary}%
        \usebeamerfont{footline}\WeekLabel\ \textperiodcentered\ \WeekTopic%
        \hfill \insertframenumber\,/\,\inserttotalframenumber%
    \end{beamercolorbox}%
}

% Halaman judul: judul, subjudul, lalu blok identitas mata kuliah
\setbeamertemplate{title page}{%
    \begin{centering}
        \vskip1.0em
        {\usebeamerfont{title}\usebeamercolor[fg]{title}\inserttitle\par}
        \vskip0.3em
        {\usebeamerfont{subtitle}\usebeamercolor[fg]{subtitle}\insertsubtitle\par}
        \vskip1.2em
        \textcolor{ubbblue}{\rule{0.72\textwidth}{0.8pt}}
        \vskip1.2em
        {\small
            \textbf{\CourseName}\ (\CourseCode)\par
            \vskip0.4em
            \CourseProgram\par
            \CourseFaculty\par
            \CourseInstitution\par
            \vskip0.6em
            \CourseSemester\ \textperiodcentered\ \WeekLabel\par
            \vskip0.8em
            Dosen Pengampu: \CourseLecturer\par
        }
        \vskip1.0em
        \textcolor{ubbblue}{\rule{0.72\textwidth}{0.4pt}}
    \end{centering}
}

% =====================================================================
%  PERINTAH BANTU
% =====================================================================
\newcommand{\R}{\mathbb{R}}
\newcommand{\N}{\mathbb{N}}
\newcommand{\Z}{\mathbb{Z}}
\newcommand{\code}[1]{\texttt{#1}}

% Rujukan gambar dari buku: \bookfig[lebar]{NN}{namaberkas.png}
% Jalur relatif terhadap folder ppt/ (dikompilasi dari sana).
\newcommand{\bookfig}[3][0.80\linewidth]{%
    \includegraphics[width=#1]{../buku_ajar/chapters/#2/figures/#3}%
}

% Keterangan sumber gambar, sepadan dengan \sumber di buku
\newcommand{\figsumber}[1]{%
    \par\vspace{1pt}{\tiny\itshape Sumber: #1\par}%
}

% Penanda blok materi yang sering dipakai di slide
\newcommand{\slidejudul}[1]{{\large\bfseries\textcolor{ubbblue}{#1}}\par\vspace{2pt}}
\newcommand{\kunci}[1]{\textcolor{ubbred}{\textbf{#1}}}

%      \title{...} \subtitle{...} \setWeekInfo{Minggu n}{Topik}
%
%  Berkas minggu hanya berisi \setWeekInfo dan lingkungan frame.
%  Kompilasi dijalankan DARI DALAM folder ppt/ (lihat compile-ppt.bat),
%  karena gambar buku dirujuk dengan jalur relatif ../buku_ajar/.
% =====================================================================

% ---------------------------------------------------------------- bahasa
\usepackage[utf8]{inputenc}
\usepackage[T1]{fontenc}
\usepackage[indonesian]{babel}
\usepackage{lmodern}
\usepackage{microtype}

% ------------------------------------------------------------- matematika
\usepackage{amsmath,amssymb,mathtools}

% ----------------------------------------------------------------- tabel
\usepackage{booktabs}
\usepackage{array}
\usepackage{tabularx}

% --------------------------------------------------------------- grafik
\usepackage{tikz}
\usetikzlibrary{arrows.meta,positioning,shapes.geometric,calc,fit,backgrounds}

% ---------------------------------------------------------- kode program
\usepackage{listings}
\lstdefinestyle{slidecode}{
    basicstyle=\ttfamily\scriptsize,
    keywordstyle=\color{ubbblue}\bfseries,
    commentstyle=\color{ubbgreen}\itshape,
    stringstyle=\color{ubbred},
    numbers=left,
    numberstyle=\tiny\color{gray},
    numbersep=6pt,
    frame=single,
    rulecolor=\color{ubbblue!40},
    backgroundcolor=\color{ubbblue!3},
    breaklines=true,
    showstringspaces=false,
    tabsize=4,
    columns=fullflexible,
    keepspaces=true,
}
\lstset{style=slidecode}

% =====================================================================
%  WARNA INSTITUSI
% =====================================================================
\definecolor{ubbblue}{RGB}{12, 62, 120}      % warna utama
\definecolor{ubbteal}{RGB}{0, 118, 110}      % aksen kedua
\definecolor{ubbgreen}{RGB}{22, 122, 74}     % untuk blok contoh
\definecolor{ubbred}{RGB}{168, 42, 42}       % untuk blok peringatan
\definecolor{ubbgray}{RGB}{70, 76, 84}       % teks sekunder

% =====================================================================
%  IDENTITAS MATA KULIAH
% =====================================================================
\newcommand{\CourseName}{Pengantar Kriptografi}
\newcommand{\CourseCode}{MPFA014}
\newcommand{\CourseSemester}{Semester Ganjil}
\newcommand{\CourseProgram}{Program Studi Matematika}
\newcommand{\CourseFaculty}{Fakultas Matematika dan Ilmu Pengetahuan Alam}
\newcommand{\CourseInstitution}{Universitas Bale Bandung}
\newcommand{\CourseLecturer}{Cecep Suwanda, S.Si., M.Kom.}

% Label minggu, disetel tiap berkas lewat \setWeekInfo
\newcommand{\WeekLabel}{}
\newcommand{\WeekTopic}{}
\newcommand{\setWeekInfo}[2]{%
    \renewcommand{\WeekLabel}{#1}%
    \renewcommand{\WeekTopic}{#2}%
}

% =====================================================================
%  TEMA: header, footer, warna blok
% =====================================================================
\usetheme{default}
\setbeamertemplate{navigation symbols}{}
\setbeamertemplate{blocks}[rounded][shadow=false]

\setbeamercolor{palette primary}{bg=ubbblue, fg=white}
\setbeamercolor{palette secondary}{bg=ubbgray, fg=white}
\setbeamercolor{palette tertiary}{bg=ubbteal, fg=white}

\setbeamercolor{structure}{fg=ubbblue}
\setbeamercolor{normal text}{fg=black!92}
\setbeamercolor{alerted text}{fg=ubbred}

\setbeamercolor{block title}{bg=ubbblue, fg=white}
\setbeamercolor{block body}{bg=ubbblue!6, fg=black!92}
\setbeamercolor{block title alerted}{bg=ubbred, fg=white}
\setbeamercolor{block body alerted}{bg=ubbred!6, fg=black!92}
\setbeamercolor{block title example}{bg=ubbgreen, fg=white}
\setbeamercolor{block body example}{bg=ubbgreen!6, fg=black!92}

\setbeamercolor{itemize item}{fg=ubbblue}
\setbeamercolor{itemize subitem}{fg=ubbteal}
\setbeamercolor{enumerate item}{fg=ubbblue}

\setbeamerfont{headline}{size=\scriptsize}
\setbeamerfont{footline}{size=\scriptsize}
\setbeamerfont{block title}{size=\normalsize}

% Badan blok rata kiri (bukan rata kanan-kiri): teks slide lebih mudah dibaca
% dan baris terakhir yang pendek tidak lagi memicu "Underfull \hbox".
\addtobeamertemplate{block body begin}{\raggedright}{}
\addtobeamertemplate{block alerted body begin}{\raggedright}{}
\addtobeamertemplate{block example body begin}{\raggedright}{}

% Header: nama mata kuliah dan kode MK
\setbeamertemplate{headline}{%
    \begin{beamercolorbox}[wd=\paperwidth,ht=2.6ex,dp=1.2ex,leftskip=1em,rightskip=1em]{palette primary}%
        \usebeamerfont{headline}\CourseName\ \textperiodcentered\ \CourseCode%
        \hfill \CourseProgram%
    \end{beamercolorbox}%
}

% Footer: minggu dan nomor halaman
\setbeamertemplate{footline}{%
    \begin{beamercolorbox}[wd=\paperwidth,ht=2.4ex,dp=1.2ex,leftskip=1em,rightskip=1em]{palette secondary}%
        \usebeamerfont{footline}\WeekLabel\ \textperiodcentered\ \WeekTopic%
        \hfill \insertframenumber\,/\,\inserttotalframenumber%
    \end{beamercolorbox}%
}

% Halaman judul: judul, subjudul, lalu blok identitas mata kuliah
\setbeamertemplate{title page}{%
    \begin{centering}
        \vskip1.0em
        {\usebeamerfont{title}\usebeamercolor[fg]{title}\inserttitle\par}
        \vskip0.3em
        {\usebeamerfont{subtitle}\usebeamercolor[fg]{subtitle}\insertsubtitle\par}
        \vskip1.2em
        \textcolor{ubbblue}{\rule{0.72\textwidth}{0.8pt}}
        \vskip1.2em
        {\small
            \textbf{\CourseName}\ (\CourseCode)\par
            \vskip0.4em
            \CourseProgram\par
            \CourseFaculty\par
            \CourseInstitution\par
            \vskip0.6em
            \CourseSemester\ \textperiodcentered\ \WeekLabel\par
            \vskip0.8em
            Dosen Pengampu: \CourseLecturer\par
        }
        \vskip1.0em
        \textcolor{ubbblue}{\rule{0.72\textwidth}{0.4pt}}
    \end{centering}
}

% =====================================================================
%  PERINTAH BANTU
% =====================================================================
\newcommand{\R}{\mathbb{R}}
\newcommand{\N}{\mathbb{N}}
\newcommand{\Z}{\mathbb{Z}}
\newcommand{\code}[1]{\texttt{#1}}

% Rujukan gambar dari buku: \bookfig[lebar]{NN}{namaberkas.png}
% Jalur relatif terhadap folder ppt/ (dikompilasi dari sana).
\newcommand{\bookfig}[3][0.80\linewidth]{%
    \includegraphics[width=#1]{../buku_ajar/chapters/#2/figures/#3}%
}

% Keterangan sumber gambar, sepadan dengan \sumber di buku
\newcommand{\figsumber}[1]{%
    \par\vspace{1pt}{\tiny\itshape Sumber: #1\par}%
}

% Penanda blok materi yang sering dipakai di slide
\newcommand{\slidejudul}[1]{{\large\bfseries\textcolor{ubbblue}{#1}}\par\vspace{2pt}}
\newcommand{\kunci}[1]{\textcolor{ubbred}{\textbf{#1}}}

%      \title{...} \subtitle{...} \setWeekInfo{Minggu n}{Topik}
%
%  Berkas minggu hanya berisi \setWeekInfo dan lingkungan frame.
%  Kompilasi dijalankan DARI DALAM folder ppt/ (lihat compile-ppt.bat),
%  karena gambar buku dirujuk dengan jalur relatif ../buku_ajar/.
% =====================================================================

% ---------------------------------------------------------------- bahasa
\usepackage[utf8]{inputenc}
\usepackage[T1]{fontenc}
\usepackage[indonesian]{babel}
\usepackage{lmodern}
\usepackage{microtype}

% ------------------------------------------------------------- matematika
\usepackage{amsmath,amssymb,mathtools}

% ----------------------------------------------------------------- tabel
\usepackage{booktabs}
\usepackage{array}
\usepackage{tabularx}

% --------------------------------------------------------------- grafik
\usepackage{tikz}
\usetikzlibrary{arrows.meta,positioning,shapes.geometric,calc,fit,backgrounds}

% ---------------------------------------------------------- kode program
\usepackage{listings}
\lstdefinestyle{slidecode}{
    basicstyle=\ttfamily\scriptsize,
    keywordstyle=\color{ubbblue}\bfseries,
    commentstyle=\color{ubbgreen}\itshape,
    stringstyle=\color{ubbred},
    numbers=left,
    numberstyle=\tiny\color{gray},
    numbersep=6pt,
    frame=single,
    rulecolor=\color{ubbblue!40},
    backgroundcolor=\color{ubbblue!3},
    breaklines=true,
    showstringspaces=false,
    tabsize=4,
    columns=fullflexible,
    keepspaces=true,
}
\lstset{style=slidecode}

% =====================================================================
%  WARNA INSTITUSI
% =====================================================================
\definecolor{ubbblue}{RGB}{12, 62, 120}      % warna utama
\definecolor{ubbteal}{RGB}{0, 118, 110}      % aksen kedua
\definecolor{ubbgreen}{RGB}{22, 122, 74}     % untuk blok contoh
\definecolor{ubbred}{RGB}{168, 42, 42}       % untuk blok peringatan
\definecolor{ubbgray}{RGB}{70, 76, 84}       % teks sekunder

% =====================================================================
%  IDENTITAS MATA KULIAH
% =====================================================================
\newcommand{\CourseName}{Pengantar Kriptografi}
\newcommand{\CourseCode}{MPFA014}
\newcommand{\CourseSemester}{Semester Ganjil}
\newcommand{\CourseProgram}{Program Studi Matematika}
\newcommand{\CourseFaculty}{Fakultas Matematika dan Ilmu Pengetahuan Alam}
\newcommand{\CourseInstitution}{Universitas Bale Bandung}
\newcommand{\CourseLecturer}{Cecep Suwanda, S.Si., M.Kom.}

% Label minggu, disetel tiap berkas lewat \setWeekInfo
\newcommand{\WeekLabel}{}
\newcommand{\WeekTopic}{}
\newcommand{\setWeekInfo}[2]{%
    \renewcommand{\WeekLabel}{#1}%
    \renewcommand{\WeekTopic}{#2}%
}

% =====================================================================
%  TEMA: header, footer, warna blok
% =====================================================================
\usetheme{default}
\setbeamertemplate{navigation symbols}{}
\setbeamertemplate{blocks}[rounded][shadow=false]

\setbeamercolor{palette primary}{bg=ubbblue, fg=white}
\setbeamercolor{palette secondary}{bg=ubbgray, fg=white}
\setbeamercolor{palette tertiary}{bg=ubbteal, fg=white}

\setbeamercolor{structure}{fg=ubbblue}
\setbeamercolor{normal text}{fg=black!92}
\setbeamercolor{alerted text}{fg=ubbred}

\setbeamercolor{block title}{bg=ubbblue, fg=white}
\setbeamercolor{block body}{bg=ubbblue!6, fg=black!92}
\setbeamercolor{block title alerted}{bg=ubbred, fg=white}
\setbeamercolor{block body alerted}{bg=ubbred!6, fg=black!92}
\setbeamercolor{block title example}{bg=ubbgreen, fg=white}
\setbeamercolor{block body example}{bg=ubbgreen!6, fg=black!92}

\setbeamercolor{itemize item}{fg=ubbblue}
\setbeamercolor{itemize subitem}{fg=ubbteal}
\setbeamercolor{enumerate item}{fg=ubbblue}

\setbeamerfont{headline}{size=\scriptsize}
\setbeamerfont{footline}{size=\scriptsize}
\setbeamerfont{block title}{size=\normalsize}

% Badan blok rata kiri (bukan rata kanan-kiri): teks slide lebih mudah dibaca
% dan baris terakhir yang pendek tidak lagi memicu "Underfull \hbox".
\addtobeamertemplate{block body begin}{\raggedright}{}
\addtobeamertemplate{block alerted body begin}{\raggedright}{}
\addtobeamertemplate{block example body begin}{\raggedright}{}

% Header: nama mata kuliah dan kode MK
\setbeamertemplate{headline}{%
    \begin{beamercolorbox}[wd=\paperwidth,ht=2.6ex,dp=1.2ex,leftskip=1em,rightskip=1em]{palette primary}%
        \usebeamerfont{headline}\CourseName\ \textperiodcentered\ \CourseCode%
        \hfill \CourseProgram%
    \end{beamercolorbox}%
}

% Footer: minggu dan nomor halaman
\setbeamertemplate{footline}{%
    \begin{beamercolorbox}[wd=\paperwidth,ht=2.4ex,dp=1.2ex,leftskip=1em,rightskip=1em]{palette secondary}%
        \usebeamerfont{footline}\WeekLabel\ \textperiodcentered\ \WeekTopic%
        \hfill \insertframenumber\,/\,\inserttotalframenumber%
    \end{beamercolorbox}%
}

% Halaman judul: judul, subjudul, lalu blok identitas mata kuliah
\setbeamertemplate{title page}{%
    \begin{centering}
        \vskip1.0em
        {\usebeamerfont{title}\usebeamercolor[fg]{title}\inserttitle\par}
        \vskip0.3em
        {\usebeamerfont{subtitle}\usebeamercolor[fg]{subtitle}\insertsubtitle\par}
        \vskip1.2em
        \textcolor{ubbblue}{\rule{0.72\textwidth}{0.8pt}}
        \vskip1.2em
        {\small
            \textbf{\CourseName}\ (\CourseCode)\par
            \vskip0.4em
            \CourseProgram\par
            \CourseFaculty\par
            \CourseInstitution\par
            \vskip0.6em
            \CourseSemester\ \textperiodcentered\ \WeekLabel\par
            \vskip0.8em
            Dosen Pengampu: \CourseLecturer\par
        }
        \vskip1.0em
        \textcolor{ubbblue}{\rule{0.72\textwidth}{0.4pt}}
    \end{centering}
}

% =====================================================================
%  PERINTAH BANTU
% =====================================================================
\newcommand{\R}{\mathbb{R}}
\newcommand{\N}{\mathbb{N}}
\newcommand{\Z}{\mathbb{Z}}
\newcommand{\code}[1]{\texttt{#1}}

% Rujukan gambar dari buku: \bookfig[lebar]{NN}{namaberkas.png}
% Jalur relatif terhadap folder ppt/ (dikompilasi dari sana).
\newcommand{\bookfig}[3][0.80\linewidth]{%
    \includegraphics[width=#1]{../buku_ajar/chapters/#2/figures/#3}%
}

% Keterangan sumber gambar, sepadan dengan \sumber di buku
\newcommand{\figsumber}[1]{%
    \par\vspace{1pt}{\tiny\itshape Sumber: #1\par}%
}

% Penanda blok materi yang sering dipakai di slide
\newcommand{\slidejudul}[1]{{\large\bfseries\textcolor{ubbblue}{#1}}\par\vspace{2pt}}
\newcommand{\kunci}[1]{\textcolor{ubbred}{\textbf{#1}}}


\title{Kriptografi Klasik II: Abjad-Majemuk dan Polygram}
\subtitle{Vigen\`ere, Kasiski, Playfair, Hill, dan Enigma}
\setWeekInfo{Minggu 5}{Kriptografi klasik II}

\begin{document}

% ----------------------------------------------------------- halaman judul
\begin{frame}[plain]
    \titlepage
\end{frame}

% ------------------------------------------------------------------ peta
\begin{frame}{Peta Materi Pertemuan Ini}
    \begin{columns}[T,onlytextwidth]
        \begin{column}{0.52\textwidth}
            \slidejudul{Lima bagian}
            \begin{enumerate}
                \item Analisis frekuensi dan batasnya
                \item Vigen\`ere dan metode Kasiski
                \item Playfair: substitusi bigram
                \item Hill: cipher blok aljabar linier
                \item Enigma: cipher berbasis mesin
            \end{enumerate}
        \end{column}
        \begin{column}{0.44\textwidth}
            \begin{block}{Referensi}
                \footnotesize
                Bab 4 buku ajar; Munir, \textit{IF4020 Kriptografi},
                ITB, 2026 (modul 04)
            \end{block}
            \begin{exampleblock}{Benang merah}
                \footnotesize
                Setiap cipher memperbesar unit yang disubstitusi, dan setiap
                pembesaran itu memunculkan kelemahan baru yang khas.
            \end{exampleblock}
        \end{column}
    \end{columns}
\end{frame}

% ---------------------------------------------------- capaian pembelajaran
\begin{frame}{Capaian Pembelajaran}
    \begin{block}{Sub-CPMK 1.1}
        \footnotesize
        Mengenal cipher klasik, yaitu substitusi (Caesar, monoalfabetik,
        Affine), transposisi (Rail Fence), Vigen\`{e}re, Playfair, Hill,
        dan Enigma, beserta gagasan kriptanalisisnya (analisis frekuensi
        dan metode Kasiski).
    \end{block}
    \vspace{0.12cm}
    \begin{exampleblock}{Pada pertemuan ini, mahasiswa mampu}
        \footnotesize
        \begin{itemize}
            \item menjelaskan mengapa analisis frekuensi melemah pada cipher
                  abjad-majemuk;
            \item menerapkan rumus enkripsi dan dekripsi Vigen\`ere;
            \item menjelaskan langkah metode Kasiski untuk menduga panjang
                  kunci;
            \item menerapkan tiga aturan Playfair pada bigram;
            \item menghitung enkripsi dan dekripsi Hill;
            \item menjelaskan prinsip rotor dan kriptanalisis Enigma.
        \end{itemize}
    \end{exampleblock}
\end{frame}

% =====================================================================
\section{Analisis Frekuensi dan Batasnya}
% =====================================================================

\begin{frame}{Analisis Frekuensi: Gagasan}
    \begin{columns}[T,onlytextwidth]
        \begin{column}{0.56\textwidth}
            \footnotesize
            \begin{itemize}
                \item Cipher abjad-tunggal memetakan satu huruf ke satu huruf
                      cipherteks yang tetap.
                \item Akibatnya sebaran frekuensi plainteks tercermin pada
                      cipherteks.
                \item Huruf tersering bahasa Inggris: \code{E}, \code{T},
                      \code{A}, \code{O}, \code{I}, \code{N}, \code{S},
                      \code{H}, \code{R}, \code{D}, \code{L}, \code{U}.
                \item Bigram tersering: \code{TH}, \code{HE}, \code{IN},
                      \code{EN}, \code{NT}, \code{RE}, \code{ER}, \code{AN},
                      \code{TI}, \code{ES}.
                \item Trigram tersering: \code{THE}, \code{AND}, \code{THA},
                      \code{ENT}, \code{ING}, \code{ION}, \code{TIO},
                      \code{FOR}, \code{NDE}, \code{HAS}.
            \end{itemize}
        \end{column}
        \begin{column}{0.42\textwidth}
            \centering
            \bookfig[0.95\linewidth]{04}{frekuensi-huruf-inggris.png}
            \figsumber{Munir, \textit{IF4020 Kriptografi}, ITB.}
        \end{column}
    \end{columns}
\end{frame}

\begin{frame}{Analisis Frekuensi: Langkah Kerja}
    \begin{columns}[T,onlytextwidth]
        \begin{column}{0.54\textwidth}
            \footnotesize
            \begin{enumerate}
                \item Hitung frekuensi relatif setiap huruf cipherteks.
                \item Bandingkan dengan tabel bahasa; huruf tersering
                      cipherteks diduga \code{E} pada plainteks.
                \item Ulangi untuk huruf tersering berikutnya; biasanya dua
                      sampai tiga huruf sudah cukup.
                \item Lanjutkan dengan bigram, trigram, dan seterusnya.
            \end{enumerate}
            \begin{alertblock}{Yang dicari bukan kunci}
                \footnotesize
                Kriptanalis menebak pemetaan huruf, bukan menelusuri ruang
                kunci.
            \end{alertblock}
        \end{column}
        \begin{column}{0.44\textwidth}
            \centering
            \bookfig[0.95\linewidth]{04}{tabel-frekuensi-huruf.png}
            \figsumber{Munir, \textit{IF4020 Kriptografi}, ITB.}
        \end{column}
    \end{columns}
\end{frame}

\begin{frame}{Mengapa Perlu Cipher Abjad-Majemuk}
    \footnotesize
    \begin{itemize}
        \item Analisis frekuensi bekerja karena satu huruf selalu menjadi
              huruf yang sama.
        \item Pertahanan alaminya: buat satu huruf plainteks menjadi
              huruf cipherteks yang berbeda-beda.
        \item Cipher abjad-majemuk memakai kunci berbeda untuk setiap posisi.
        \item Bila kunci berulang secara periodik, cipher menjadi rangkaian
              Caesar dengan pergeseran berganti-ganti.
        \item Akibatnya histogram cipherteks menjadi lebih datar.
    \end{itemize}
    \vspace{0.15cm}
    \begin{block}{Harga yang dibayar}
        \footnotesize
        Pola statistik huruf tunggal tertutup, tetapi periodisitas kunci
        meninggalkan jejak yang dapat dilacak.
    \end{block}
\end{frame}

% =====================================================================
\section{Vigen\`ere dan Metode Kasiski}
% =====================================================================

\begin{frame}{Vigen\`ere: Kunci yang Diulang}
    \begin{columns}[T,onlytextwidth]
        \begin{column}{0.54\textwidth}
            \footnotesize
            Kunci $K=k_1k_2\cdots k_m$ berpanjang $m$ diulang secara periodik
            sepanjang plainteks.
            \[
                c_j=(p_j+k_i)\bmod 26,
                \qquad
                p_j=(c_j-k_i)\bmod 26.
            \]
            \begin{itemize}
                \item Setiap huruf dienkripsi dengan Caesar memakai huruf kunci
                      pada posisi itu.
                \item Untuk $m=1$, Vigen\`ere kembali menjadi cipher
                      abjad-tunggal.
                \item Tabel Vigen\`ere: baris ke-$i$ adalah alfabet yang
                      digeser $i$ posisi.
            \end{itemize}
        \end{column}
        \begin{column}{0.42\textwidth}
            \centering
            \bookfig[1.0\linewidth]{04}{vigenere-square.png}
            \figsumber{Munir, \textit{IF4020 Kriptografi}, ITB.}
        \end{column}
    \end{columns}
\end{frame}

\begin{frame}{Vigen\`ere: Contoh Tiga Huruf Pertama}
    \footnotesize
    Kunci \code{soreang} bernilai $(18,14,17,4,0,13,6)$. Plainteks
    \code{terimapesanangulaikambing} berisi 25 huruf, sehingga kunci diulang
    menjadi \code{soreangsoreangsoreangsore}.
    \vspace{0.08cm}
    \begin{center}
    \begin{tabular}{@{}lccc@{}}
        \toprule
         & Huruf ke-1 & Huruf ke-2 & Huruf ke-3 \\
        \midrule
        Plainteks $P$ & $t=19$ & $e=4$ & $r=17$ \\
        Kunci $K$ & $s=18$ & $o=14$ & $r=17$ \\
        $C=(P+K)\bmod 26$ & $(19+18)\bmod 26=11$ &
        $(4+14)\bmod 26=18$ & $(17+17)\bmod 26=8$ \\
        Cipherteks & \code{L} & \code{S} & \code{I} \\
        \bottomrule
    \end{tabular}
    \end{center}
    \vspace{0.08cm}
    \begin{exampleblock}{Pemeriksaan dekripsi}
        \footnotesize
        $L=11$ dikurangi $s=18$ memberi
        $(11-18)\bmod 26=-7\bmod 26=19=t$.
    \end{exampleblock}
\end{frame}

\begin{frame}{Vigen\`ere: Hasil dan Sifatnya}
    \begin{columns}[T,onlytextwidth]
        \begin{column}{0.56\textwidth}
            \footnotesize
            \begin{itemize}
                \item Melanjutkan perhitungan yang sama memberi cipherteks
                      \code{LSIMMNVWGRRAAMMZRMKNSTWEK}.
                \item Huruf \code{e} menjadi \code{S} pada posisi kedua dan
                      \code{R} pada posisi kedelapan.
                \item Huruf cipherteks \code{M} dapat berasal dari \code{i},
                      \code{m}, \code{g}, \code{u}, atau \code{k}.
                \item Distribusi huruf cipherteks menjadi lebih seragam,
                      sehingga histogramnya cenderung datar.
            \end{itemize}
            \begin{alertblock}{Akibatnya}
                \footnotesize
                Analisis frekuensi huruf tunggal kehilangan pijakan.
            \end{alertblock}
        \end{column}
        \begin{column}{0.42\textwidth}
            \centering
            \bookfig[1.0\linewidth]{04}{histogram-flat.png}
            \figsumber{Munir, \textit{IF4020 Kriptografi}, ITB.}
        \end{column}
    \end{columns}
\end{frame}

\begin{frame}{Vigen\`ere: Sejarah Singkat}
    \begin{columns}[T,onlytextwidth]
        \begin{column}{0.58\textwidth}
            \footnotesize
            \begin{itemize}
                \item Dijelaskan Giovan Batista Belaso pada tahun 1553.
                \item Dipublikasikan Blaise de Vigen\`ere pada tahun 1586.
                \item Pada abad ke-19 banyak orang mengira Vigen\`ere
                      penemunya, sehingga namanya melekat.
                \item Dipecahkan Babbage dan Kasiski pada pertengahan abad
                      ke-19.
                \item Dipakai Tentara Konfederasi pada Perang Sipil Amerika.
            \end{itemize}
        \end{column}
        \begin{column}{0.38\textwidth}
            \centering
            \bookfig[0.72\linewidth]{04}{blaise-de-vigenere.png}\\
            {\scriptsize Blaise de Vigen\`ere (1523--1596)}
            \figsumber{Munir, \textit{IF4020 Kriptografi}, ITB.}
        \end{column}
    \end{columns}
\end{frame}

\begin{frame}{Vigen\`ere: Tiga Varian}
    \footnotesize
    \begin{itemize}
        \item \textbf{Full Vigen\`ere} --- setiap baris tabel berupa permutasi
              huruf, bukan pergeseran; tabelnya dirahasiakan.
        \item \textbf{Auto-Key} --- kunci yang pendek disambung dengan
              plainteks.
        \item \textbf{Running-Key} --- kunci panjang diambil dari teks
              bermakna.
    \end{itemize}
    \vspace{0.1cm}
    \begin{exampleblock}{Auto-Key dengan kunci \code{INDO}}
        \scriptsize
        \begin{tabular}{@{}l l@{}}
            Plainteks  & \code{negarapenghasilminyakmentahdidunia}\\
            Kunci      & \code{INDONEGARAPENGHASILMINYAKMENTAHDID}\\
            Cipherteks & \code{VRJOEEVEEGWEFOSMAVJMSZCNDMLQBDBQQD}
        \end{tabular}
    \end{exampleblock}
\end{frame}

\begin{frame}{Metode Kasiski: Prinsip}
    \begin{columns}[T,onlytextwidth]
        \begin{column}{0.54\textwidth}
            \footnotesize
            \begin{itemize}
                \item Ditemukan Friedrich Kasiski pada tahun 1863.
                \item Jika sebuah string plainteks muncul berulang pada jarak
                      kelipatan panjang kunci, kedua kemunculan menghasilkan
                      kriptogram yang sama.
                \item Contoh: plainteks \code{cryptoisshortforcryptography}
                      dengan kunci \code{abcd} berpanjang 4.
                \item Kata \code{crypto} muncul dua kali pada jarak 16 huruf.
                \item Karena 16 kelipatan 4, kedua kemunculan menghasilkan
                      kriptogram \code{CSASTP}.
                \item Dengan kunci \code{abcdef} berpanjang 6, jarak 16 bukan
                      kelipatan 6, sehingga kriptogramnya berbeda.
            \end{itemize}
        \end{column}
        \begin{column}{0.42\textwidth}
            \centering
            \bookfig[1.0\linewidth]{04}{kasiski-jarak16.png}
            \figsumber{Munir, \textit{IF4020 Kriptografi}, ITB.}
        \end{column}
    \end{columns}
\end{frame}

\begin{frame}{Metode Kasiski: Empat Langkah}
    \footnotesize
    \begin{enumerate}
        \item Temukan semua kriptogram yang berulang di dalam cipherteks.
        \item Hitung jarak antara kemunculan-kemunculan itu.
        \item Faktorkan setiap jarak; faktor-faktornya menyatakan panjang
              kunci yang mungkin.
        \item Ambil irisan himpunan faktor; nilai yang muncul di irisan
              adalah dugaan panjang kunci.
    \end{enumerate}
    \vspace{0.2cm}
    \begin{alertblock}{Batas metode}
        \footnotesize
        Kasiski hanya menduga \emph{panjang} kunci, belum menemukan
        huruf-huruf kuncinya.
    \end{alertblock}
\end{frame}

\begin{frame}{Kasiski: Contoh Pendek}
    \footnotesize
    \begin{itemize}
        \item Kriptogram \code{DYDUXRMH} berulang pada jarak 18, dengan faktor
              $\{18,9,6,3,2\}$.
        \item Kriptogram \code{NQD} berulang pada jarak 20, dengan faktor
              $\{20,10,5,4,2\}$.
        \item Irisan kedua himpunan faktor adalah $\{2\}$.
        \item Jadi panjang kunci kemungkinan besar adalah 2.
    \end{itemize}
    \vspace{0.15cm}
    \begin{block}{Membaca hasilnya}
        \footnotesize
        Setiap jarak menyumbang satu himpunan kandidat; kandidat yang bertahan
        di semua himpunan adalah panjang kunci yang paling mungkin.
    \end{block}
\end{frame}

\begin{frame}{Kasiski: Contoh Lengkap \code{LJV}}
    \begin{columns}[T,onlytextwidth]
        \begin{column}{0.50\textwidth}
            \footnotesize
            \begin{itemize}
                \item Kriptogram \code{LJV} muncul enam kali.
                \item Jarak berturut-turut: 15, 15, 15, 10, dan 10.
                \item Faktor 15: $\{3,5,15\}$.
                \item Faktor 10: $\{2,5,10\}$.
                \item Satu-satunya nilai yang muncul di kedua himpunan adalah
                      5.
                \item Jadi panjang kunci adalah 5.
            \end{itemize}
        \end{column}
        \begin{column}{0.46\textwidth}
            \centering
            \bookfig[1.0\linewidth]{04}{kasiski-ljv.png}
            \figsumber{Munir, \textit{IF4020 Kriptografi}, ITB.}
        \end{column}
    \end{columns}
\end{frame}

\begin{frame}{Kasiski: Dari \code{LJV} ke Kunci \code{SCRAM}}
    \begin{columns}[T,onlytextwidth]
        \begin{column}{0.54\textwidth}
            \footnotesize
            \begin{itemize}
                \item Dengan panjang kunci 5, cipherteks dibagi menjadi lima
                      kelompok.
                \item Huruf tersering tiap kelompok: \code{L}, \code{J},
                      \code{V}, \code{N}, dan \code{A}.
                \item \code{LJV} diduga padanan trigram \code{THE}.
                \item Setiap huruf kunci dihitung dengan
                      $K=(C-P)\bmod 26$.
            \end{itemize}
            \begin{center}
            \scriptsize
            \begin{tabular}{@{}lccccc@{}}
                \toprule
                 & 1 & 2 & 3 & 4 & 5 \\
                \midrule
                Cipher & $L=11$ & $J=9$ & $V=21$ & $N=13$ & $A=0$ \\
                Plain  & $T=19$ & $H=7$ & $E=4$ & $N=13$ & $O=14$ \\
                Kunci  & $18=S$ & $2=C$ & $17=R$ & $0=A$ & $12=M$ \\
                \bottomrule
            \end{tabular}
            \end{center}
            {\footnotesize Kunci yang diperoleh: \code{SCRAM}.}
        \end{column}
        \begin{column}{0.42\textwidth}
            \centering
            \bookfig[1.0\linewidth]{04}{kasiski-scram.png}
            \figsumber{Munir, \textit{IF4020 Kriptografi}, ITB.}
        \end{column}
    \end{columns}
\end{frame}

% =====================================================================
\section{Playfair}
% =====================================================================

\begin{frame}{Playfair: Substitusi Bigram}
    \begin{columns}[T,onlytextwidth]
        \begin{column}{0.56\textwidth}
            \footnotesize
            \begin{itemize}
                \item Ditemukan Sir Charles Wheatstone dan dipromosikan Baron
                      Lyon Playfair pada tahun 1854.
                \item Satu blok plainteks disubstitusi dengan satu blok
                      cipherteks; unitnya bigram.
                \item Kunci disusun dalam matriks $5\times5$.
                \item Karena hanya ada 25 sel, huruf \code{J} dibuang dan
                      setiap \code{J} diganti \code{I}.
                \item Jumlah matriks kunci:
                      $25!=15.511.043.330.985.984.000.000$.
            \end{itemize}
        \end{column}
        \begin{column}{0.40\textwidth}
            \centering
            \bookfig[0.90\linewidth]{04}{playfair-matriks.png}
            \figsumber{Munir, \textit{IF4020 Kriptografi}, ITB.}
        \end{column}
    \end{columns}
\end{frame}

\begin{frame}{Playfair: Dari Kalimat Kunci ke Matriks}
    \begin{columns}[T,onlytextwidth]
        \begin{column}{0.52\textwidth}
            \footnotesize
            \begin{enumerate}
                \item Pilih kalimat kunci yang mudah diingat:
                      \code{JALAN GANESHA SEPULUH}.
                \item Buang huruf berulang dan huruf \code{J}, sehingga
                      tersisa \code{ALNGESHPU}.
                \item Tambahkan huruf yang belum ada, kecuali \code{J}:
                      \code{ALNGESHPUBCDFIKMOQRTVWXYZ}.
                \item Isikan baris demi baris ke dalam matriks.
            \end{enumerate}
        \end{column}
        \begin{column}{0.44\textwidth}
            \centering
            \footnotesize
            \begin{tabular}{|c|c|c|c|c|}
                \hline
                A & L & N & G & E \\ \hline
                S & H & P & U & B \\ \hline
                C & D & F & I & K \\ \hline
                M & O & Q & R & T \\ \hline
                V & W & X & Y & Z \\ \hline
            \end{tabular}
        \end{column}
    \end{columns}
\end{frame}

\begin{frame}{Playfair: Menyiapkan Plainteks}
    \footnotesize
    Contoh plainteks \code{temui ibu nanti malam}.
    \begin{enumerate}
        \item Buang spasi: \code{temuibunantimalam}.
        \item Ganti \code{j} dengan \code{i}; pada contoh ini tidak ada
              \code{j}.
        \item Tulis sebagai bigram: \code{te mu ii bu na nt im al am}.
        \item Sisipkan \code{x} pada bigram kembar:
              \code{te mu ix ib un an ti ma la m}.
        \item Bila bigram terakhir hanya satu huruf, tambahkan \code{x}:
              \code{te mu ix ib un an ti ma la mx}.
    \end{enumerate}
    \begin{alertblock}{Mengapa perlu}
        \footnotesize
        Aturan enkripsi tidak terdefinisi bila kedua huruf berada pada sel
        yang sama.
    \end{alertblock}
\end{frame}

\begin{frame}{Playfair: Tiga Aturan Enkripsi}
    \footnotesize
    \begin{tabularx}{\linewidth}{@{}>{\raggedright\arraybackslash}p{3.0cm}
        >{\raggedright\arraybackslash}X
        >{\raggedright\arraybackslash}p{4.4cm}@{}}
        \toprule
        Posisi kedua huruf & Aturan & Contoh \\
        \midrule
        Baris yang sama & Tiap huruf diganti huruf di kanannya; di ujung baris
        siklik kembali ke awal. & \code{di} $\rightarrow$ \code{FK};
        \code{qt} $\rightarrow$ \code{RM} \\
        \addlinespace
        Kolom yang sama & Tiap huruf diganti huruf di bawahnya; di ujung kolom
        siklik kembali ke awal. & \code{nq} $\rightarrow$ \code{PX};
        \code{ow} $\rightarrow$ \code{WL} \\
        \addlinespace
        Membentuk persegi panjang & Huruf pertama mengambil barisnya sendiri
        dan kolom huruf kedua; huruf kedua berlaku sebaliknya. &
        \code{hz} $\rightarrow$ \code{BW};
        \code{IH} $\rightarrow$ \code{DU} \\
        \bottomrule
    \end{tabularx}
    \vspace{0.15cm}
    \begin{block}{Dekripsi}
        \footnotesize
        Baris yang sama: huruf di kirinya. Kolom yang sama: huruf di atasnya.
        Aturan persegi panjang berlaku sama. Huruf pengisi \code{X} kemudian
        dibuang.
    \end{block}
\end{frame}

\begin{frame}{Playfair: Contoh Terhitung}
    \begin{columns}[T,onlytextwidth]
        \begin{column}{0.52\textwidth}
            \footnotesize
            Bigram yang telah disiapkan:
            \begin{center}
            \scriptsize\texttt{te mu ix ib un an ti ma la mx}
            \end{center}
            \begin{itemize}
                \item Bigram pertama \code{te}: \code{t} di baris 4 kolom 5,
                      \code{e} di baris 1 kolom 5.
                \item Keduanya satu kolom, sehingga masing-masing diganti
                      huruf di bawahnya.
                \item $t\rightarrow Z$ dan $e\rightarrow B$, jadi
                      \code{te} $\rightarrow$ \code{ZB}.
                \item Melanjutkan seluruh bigram memberi
                      \code{ZB RS FY KU PG LG RK VS NL QV}.
            \end{itemize}
        \end{column}
        \begin{column}{0.44\textwidth}
            \centering
            \bookfig[1.0\linewidth]{04}{playfair-contoh.png}
            \figsumber{Munir, \textit{IF4020 Kriptografi}, ITB.}
        \end{column}
    \end{columns}
\end{frame}

\begin{frame}{Playfair: Kriptanalisis}
    \footnotesize
    \begin{itemize}
        \item Terdapat $26\times26=676$ bigram, sehingga identifikasi bigram
              jauh lebih sukar daripada huruf tunggal.
        \item Namun unitnya hanya dua huruf, sehingga cipher ini tetap tidak
              aman.
        \item Frekuensi bigram bahasa masih bocor: \code{TH} dan \code{HE}
              paling sering muncul.
        \item Bigram dan kebalikannya, misalnya \code{AB} dan \code{BA},
              menghasilkan pola plainteks yang sama; contohnya \code{RE} dan
              \code{ER}.
        \item Tebakan \code{TH} $\rightarrow$ \code{BM} memaksa posisi
              \code{T}, \code{H}, \code{B}, dan \code{M} konsisten, lalu
              matriks parsial dapat disusun.
    \end{itemize}
    \vspace{0.15cm}
    \begin{exampleblock}{Pelajaran}
        \footnotesize
        Playfair menaikkan biaya analisis frekuensi, tetapi tidak menghapus
        struktur statistik bahasa.
    \end{exampleblock}
\end{frame}

% =====================================================================
\section{Hill}
% =====================================================================

\begin{frame}{Hill: Cipher Blok Berbasis Aljabar Linier}
    \begin{columns}[T,onlytextwidth]
        \begin{column}{0.58\textwidth}
            \footnotesize
            \begin{itemize}
                \item Dikembangkan Lester Hill pada tahun 1929.
                \item $m$ huruf plainteks dienkripsi sekaligus menjadi $m$
                      huruf cipherteks.
                \item Untuk $m=3$:
                      $C_1=(k_{11}p_1+k_{12}p_2+k_{13}p_3)\bmod 26$,
                      dan seterusnya untuk $C_2$ dan $C_3$.
                \item Bentuk matriks: $\mathbf{C}=\mathbf{KP}\pmod{26}$,
                      dengan $\mathbf{K}$ berukuran $m\times m$ dan
                      $\mathbf{P}$ vektor kolom.
                \item Dekripsi memakai $\mathbf{K}^{-1}$.
            \end{itemize}
        \end{column}
        \begin{column}{0.40\textwidth}
            \centering
            \bookfig[0.95\linewidth]{04}{hill-matriks.png}
            \figsumber{Munir, \textit{IF4020 Kriptografi}, ITB.}
        \end{column}
    \end{columns}
\end{frame}

\begin{frame}{Hill: Syarat Kunci dan Invers}
    \begin{columns}[T,onlytextwidth]
        \begin{column}{0.56\textwidth}
            \footnotesize
            \begin{itemize}
                \item Dekripsi memerlukan $\mathbf{K}^{-1}$ dengan
                      $\mathbf{KK}^{-1}=\mathbf{I}$.
                \item Invers modulo 26 ada hanya bila
                      $\gcd(\det(\mathbf{K}),26)=1$.
                \item Untuk $2\times2$:
                      $\mathbf{K}^{-1}=(\det)^{-1}
                      \begin{pmatrix}d&-b\\-c&a\end{pmatrix}$.
                \item Untuk $3\times3$ dipakai matriks kofaktor dengan
                      $\det(\mathbf{K})=aA+bB+cC$.
            \end{itemize}
            \begin{alertblock}{Pada kunci $m=3$ nanti}
                \footnotesize
                $\det(\mathbf{K})=-939\equiv 23\pmod{26}$ dan
                $\gcd(23,26)=1$, sehingga inversnya ada.
            \end{alertblock}
        \end{column}
        \begin{column}{0.42\textwidth}
            \centering
            \bookfig[1.0\linewidth]{04}{hill-invers-3x3.png}
            \figsumber{Munir, \textit{IF4020 Kriptografi}, ITB.}
        \end{column}
    \end{columns}
\end{frame}

\begin{frame}{Hill: Contoh Enkripsi $m=3$}
    \footnotesize
    Matriks kunci
    $\mathbf{K}=\begin{pmatrix}17&17&5\\21&18&21\\2&2&19\end{pmatrix}$
    dan plainteks \code{paymoremoney}. Kelompok pertama \code{pay} memberi
    $\mathbf{P}=(15,0,24)^T$.
    \[
        \mathbf{KP}=
        \begin{pmatrix}
            17\cdot15+17\cdot0+5\cdot24\\
            21\cdot15+18\cdot0+21\cdot24\\
            2\cdot15+2\cdot0+19\cdot24
        \end{pmatrix}
        =
        \begin{pmatrix}375\\819\\486\end{pmatrix}
        \bmod 26
        =
        \begin{pmatrix}11\\13\\18\end{pmatrix}
        =\text{\code{LNS}}.
    \]
    Meneruskan cara yang sama untuk seluruh blok menghasilkan cipherteks
    \begin{center}
        \code{LNSHDLEWMTRW}.
    \end{center}
\end{frame}

\begin{frame}{Hill: Contoh Dekripsi dan Uji}
    \footnotesize
    Invers kunci yang sama adalah
    $\mathbf{K}^{-1}=
    \begin{pmatrix}4&9&15\\15&17&6\\24&0&17\end{pmatrix}$.
    Dekripsi \code{LNS}$=(11,13,18)^T$:
    \[
        \mathbf{K}^{-1}
        \begin{pmatrix}11\\13\\18\end{pmatrix}
        =
        \begin{pmatrix}431\\494\\570\end{pmatrix}
        \bmod 26
        =
        \begin{pmatrix}15\\0\\24\end{pmatrix}
        =\text{\code{PAY}}.
    \]
    \begin{block}{Uji $\mathbf{KK}^{-1}=\mathbf{I}$}
        \footnotesize
        Baris pertama $\mathbf{K}$ dikali kolom pertama $\mathbf{K}^{-1}$:
        $17\cdot4+17\cdot15+5\cdot24=443\equiv1\pmod{26}$.
    \end{block}
\end{frame}

\begin{frame}{Hill: Menghitung Invers $2\times2$}
    \footnotesize
    Ambil
    $\mathbf{K}=\begin{pmatrix}3&10\\15&9\end{pmatrix}$.
    Determinannya
    \[
        \det(\mathbf{K})=3\cdot9-15\cdot10=27-150=-123\equiv7\pmod{26}.
    \]
    Karena $7\cdot15=105\equiv1\pmod{26}$, maka $7^{-1}\equiv15$.
    Dengan $-10\equiv16$ dan $-15\equiv11$,
    \[
        \mathbf{K}^{-1}=15
        \begin{pmatrix}9&16\\11&3\end{pmatrix}
        =
        \begin{pmatrix}135&240\\165&45\end{pmatrix}
        \bmod 26
        =
        \begin{pmatrix}5&6\\9&19\end{pmatrix}.
    \]
    Uji balik:
    $\mathbf{KK}^{-1}=
    \begin{pmatrix}105&208\\156&261\end{pmatrix}\bmod26
    =\begin{pmatrix}1&0\\0&1\end{pmatrix}$.
\end{frame}

\begin{frame}{Hill: Kriptanalisis Known-Plaintext}
    \begin{columns}[T,onlytextwidth]
        \begin{column}{0.59\textwidth}
            \footnotesize
            Hill menyembunyikan frekuensi huruf tunggal, tetapi rumusnya
            linier. Untuk $m=2$, dua pasangan sudah cukup:
            \[
                (19,7)\rightarrow(0,23),
                \qquad
                (4,17)\rightarrow(12,6).
            \]
            Susun sebagai kolom:
            \[
                \mathbf{P}=\begin{pmatrix}19&4\\7&17\end{pmatrix},
                \quad
                \mathbf{C}=\begin{pmatrix}0&12\\23&6\end{pmatrix},
                \quad
                \mathbf{P}^{-1}=\begin{pmatrix}25&14\\5&5\end{pmatrix}.
            \]
            Dengan semua operasi modulo 26,
            \[
                \mathbf{K}=\mathbf{C}\mathbf{P}^{-1}
                =\begin{pmatrix}0&12\\23&6\end{pmatrix}
                \begin{pmatrix}25&14\\5&5\end{pmatrix}
                =\begin{pmatrix}8&8\\7&14\end{pmatrix}.
            \]
        \end{column}
        \begin{column}{0.39\textwidth}
            \centering
            \bookfig[1.0\linewidth]{04}{hill-known-plaintext.png}\\
            {\scriptsize Uji balik $(19,7)$:
            $(208,231)\bmod26=(0,23)$.}
            \figsumber{Munir, \textit{IF4020 Kriptografi}, ITB.}
        \end{column}
    \end{columns}
\end{frame}

% =====================================================================
\section{Enigma}
% =====================================================================

\begin{frame}{Enigma: Mesin dan Komponennya}
    \begin{columns}[T,onlytextwidth]
        \begin{column}{0.56\textwidth}
            \footnotesize
            \begin{itemize}
                \item Ditemukan dan dipatenkan Arthur Scherbius untuk tujuan
                      komersial, diplomatik, dan militer.
                \item Namanya berasal dari bahasa Latin \textit{enigmae} yang
                      berarti teka-teki.
                \item Dipakai tentara Nazi Jerman pada Perang Dunia II.
                \item Papan tombol mengetik plainteks; papan lampu menampilkan
                      cipherteks.
                \item Papan steker menukar pasangan huruf dan diubah setiap
                      hari.
                \item Tiga atau empat rotor, reflektor, dan jendela pengamat
                      posisi rotor.
            \end{itemize}
        \end{column}
        \begin{column}{0.40\textwidth}
            \centering
            \bookfig[0.82\linewidth]{04}{enigma-bagian.png}
            \figsumber{Munir, \textit{IF4020 Kriptografi}, ITB.}
        \end{column}
    \end{columns}
\end{frame}

\begin{frame}{Enigma: Rotor dan Periodenya}
    \begin{columns}[T,onlytextwidth]
        \begin{column}{0.54\textwidth}
            \footnotesize
            \begin{itemize}
                \item Setiap rotor melakukan substitusi abjad-tunggal.
                \item Hasil satu rotor menjadi masukan rotor berikutnya.
                \item Setiap kali satu huruf dienkripsi, rotor berputar satu
                      huruf.
                \item Setelah 26 putaran, rotor berikutnya maju satu posisi;
                      bekerja seperti odometer.
                \item Tiga rotor: $26^3=17.576$ huruf. Empat rotor:
                      $26^4=456.976$ huruf sebelum urutan berulang.
                \item Posisi awal rotor menjadi kunci Enigma.
            \end{itemize}
        \end{column}
        \begin{column}{0.42\textwidth}
            \centering
            \bookfig[1.0\linewidth]{04}{enigma-skema.png}
            \figsumber{Munir, \textit{IF4020 Kriptografi}, ITB.}
        \end{column}
    \end{columns}
\end{frame}

\begin{frame}{Enigma: Contoh Huruf demi Huruf}
    \begin{columns}[T,onlytextwidth]
        \begin{column}{0.54\textwidth}
            \footnotesize
            Model sederhana: alfabet enam huruf dan tiga rotor.
            \begin{itemize}
                \item \code{D} masuk rotor pertama, keluar \code{E}.
                \item \code{E} masuk rotor kedua, keluar \code{B}.
                \item \code{B} masuk rotor ketiga, keluar \code{E}.
                \item Jadi \code{D} dienkripsi menjadi \code{E}.
                \item Setelah rotor ketiga bergeser satu huruf, \code{D}
                      berikutnya menghasilkan \code{A}.
            \end{itemize}
            \begin{exampleblock}{Pesan sesungguhnya}
                \footnotesize
                \code{PESAWATMENDARAT} menghasilkan
                \code{GUFYLFRKQEJXVHM}; tiga huruf pertama
                \code{P}$\rightarrow$\code{G},
                \code{E}$\rightarrow$\code{Q},
                \code{S}$\rightarrow$\code{F}.
            \end{exampleblock}
        \end{column}
        \begin{column}{0.42\textwidth}
            \centering
            \bookfig[1.0\linewidth]{04}{enigma-rotor.png}
            \figsumber{Munir, \textit{IF4020 Kriptografi}, ITB.}
        \end{column}
    \end{columns}
\end{frame}

\begin{frame}{Enigma: Jalur Sinyal Lengkap}
    \footnotesize
    Pada mesin sesungguhnya, sinyal melewati papan steker, tiga rotor,
    reflektor, lalu kembali melalui ketiga rotor dan papan steker.
    \vspace{0.1cm}
    {\scriptsize
    \begin{tabularx}{\linewidth}{@{}c c >{\raggedright\arraybackslash}X c@{}}
        \toprule
        Huruf & Posisi rotor & Jalur sinyal & Keluaran \\
        \midrule
        \code{P} & \texttt{ABE} & steker \code{P} $\rightarrow$ R3 \code{W}
        $\rightarrow$ R2 \code{U} $\rightarrow$ R1 \code{A} $\rightarrow$
        reflektor \code{Y} $\rightarrow$ R1 \code{O} $\rightarrow$ R2
        \code{T} $\rightarrow$ R3 \code{G} $\rightarrow$ steker \code{G} &
        \textbf{\code{G}} \\
        \addlinespace
        \code{E} & \texttt{ABM} & steker \code{E} $\rightarrow$ R3 \code{W}
        $\rightarrow$ R2 \code{U} $\rightarrow$ R1 \code{A} $\rightarrow$
        reflektor \code{Y} $\rightarrow$ R1 \code{O} $\rightarrow$ R2
        \code{T} $\rightarrow$ R3 \code{Q} $\rightarrow$ steker \code{Q} &
        \textbf{\code{Q}} \\
        \addlinespace
        \code{S} & \texttt{ABG} & steker \code{S} $\rightarrow$ R3 \code{K}
        $\rightarrow$ R2 \code{G} $\rightarrow$ R1 \code{D} $\rightarrow$
        reflektor \code{H} $\rightarrow$ R1 \code{P} $\rightarrow$ R2
        \code{P} $\rightarrow$ R3 \code{F} $\rightarrow$ steker \code{F} &
        \textbf{\code{F}} \\
        \bottomrule
    \end{tabularx}
    }
    \vspace{0.1cm}
    \begin{exampleblock}{Hasil}
        \footnotesize
        Tiga huruf pertama \code{PES} menghasilkan \code{GQF}, bagian awal
        cipherteks \code{GUFYLFRKQEJXVHM}.
    \end{exampleblock}
\end{frame}

\begin{frame}{Enigma: Kriptanalisis}
    \begin{columns}[T,onlytextwidth]
        \begin{column}{0.56\textwidth}
            \footnotesize
            \begin{itemize}
                \item 1932: Marian Rejewski, Jerzy R\'o\.zycki, dan Henryk
                      Zygalski dari Polandia menemukan kriptanalisis Enigma.
                \item 1939: Jerman mendesain ulang Enigma dengan menambahkan
                      papan steker dan reflektor.
                \item Metode kriptanalisis yang lama tidak berlaku lagi.
                \item Dengan bantuan Polandia, Prancis, dan Inggris dibuat
                      mesin \textit{bombe}.
                \item \textit{Bombe} dirancang Alan Turing dan berhasil
                      memecahkan Enigma buatan Jerman.
                \item Keberhasilan itu diperkirakan memperpendek Perang Dunia
                      II menjadi dua tahun.
            \end{itemize}
        \end{column}
        \begin{column}{0.42\textwidth}
            \centering
            \bookfig[0.92\linewidth]{04}{enigma-bombe.png}
            \figsumber{Munir, \textit{IF4020 Kriptografi}, ITB.}
        \end{column}
    \end{columns}
\end{frame}

% =====================================================================
\section{Rangkuman}
% =====================================================================

\begin{frame}{Rangkuman}
    \footnotesize
    \begin{itemize}
        \item Analisis frekuensi bergantung pada pemetaan huruf yang tetap;
              cipher abjad-majemuk mematahkannya dengan mengganti kunci setiap
              posisi.
        \item Vigen\`ere memakai
              $c_j=(p_j+k_i)\bmod26$ dengan kunci periodik, sehingga huruf
              yang sama dapat menjadi huruf berbeda.
        \item Kasiski memfaktorkan jarak kriptogram berulang; irisan faktor
              menduga panjang kunci, lalu frekuensi per kelompok menemukan
              huruf kunci.
        \item Playfair bekerja pada bigram dengan matriks $5\times5$ dan tiga
              aturan; kuat terhadap frekuensi huruf tunggal, tetapi bocor pada
              frekuensi bigram.
        \item Hill memakai $\mathbf{C}=\mathbf{KP}\bmod26$ dengan syarat
              $\gcd(\det\mathbf{K},26)=1$; sifat liniernya membuatnya tumbang
              oleh \textit{known-plaintext}.
        \item Enigma adalah cipher abjad-majemuk berbasis rotor berperiode
              $26^3$ atau $26^4$; kriptanalisis Polandia dan \textit{bombe}
              Turing memecahkannya.
    \end{itemize}
\end{frame}

\begin{frame}{Yang Harus Diingat}
    \begin{block}{Empat butir kunci}
        \begin{itemize}
            \item Cipher abjad-majemuk menaikkan biaya analisis frekuensi,
                  tetapi periodisitas kunci membuka jalan bagi Kasiski.
            \item Playfair dan Hill memperbesar unit substitusi menjadi bigram
                  dan blok; Hill menambahkan struktur aljabar linier yang
                  justru mudah dibalik.
            \item Enigma adalah puncak cipher klasik berbasis mesin, dan
                  kerapuhannya terletak pada struktur, bukan panjang kunci.
            \item Kriptanalisis klasik selalu mencari struktur yang bertahan,
                  bukan mencoba seluruh kunci.
        \end{itemize}
    \end{block}
    \vspace{0.25cm}
    \begin{exampleblock}{Pertemuan berikutnya}
        \textbf{Minggu 6 --- One-time pad}
    \end{exampleblock}
\end{frame}

\end{document}
